\documentclass[runningheads,envcountsame,envcountsect]{llncs}

\usepackage[T1]{fontenc}
\usepackage{graphicx}
\usepackage{amsfonts}
\usepackage{latexsym,amssymb,amsmath}
\usepackage{placeins}
\usepackage{booktabs}
\usepackage{comment}
\usepackage{caption}
\usepackage{subcaption}

\usepackage[dvipsnames]{xcolor}
\usepackage{hyperref}
\renewcommand\UrlFont{\color{blue}\rmfamily}
\urlstyle{rm}

% Defer proofs to the appendix (same setup as the sar_plane paper). Statements
% stay in the main text and are repeated in the appendix before their proof.
\usepackage[bibliography=common,appendix=strip]{apxproof}
\newtheoremrep{theorem}{Theorem}[section]
\newtheoremrep{lemma}[theorem]{Lemma}

% Author comment macros.
\newcommand\JC[1]{{\color{Maroon} \textbf{JC}: #1}}
\newcommand\EK[1]{{\color{Blue} \textbf{EK}: #1}}
\newcommand\DK[1]{{\color{Green} \textbf{DK}: #1}}
\newcommand\OM[1]{{\color{Purple} \textbf{OM}: #1}}

% Notation macros.
\newcommand{\CR}{\mathrm{CR}}
\newcommand{\disc}{\mathrm{disc}}
\newcommand{\capt}{\mathrm{cap}}
\newcommand{\opt}{T^{*}}

\title{Landing-Only Linear Search for an Escaping Target: Discovery and Capture}
\titlerunning{Landing-Only Search for an Escaping Target}

% Working author list: same as the companion planar paper. Adjust as needed.
\author{Jared Coleman\inst{1} \and
Evangelos Kranakis\inst{2} \and
Danny Krizanc\inst{3} \and
Oscar Morales-Ponce\inst{4}}

\authorrunning{J. Coleman, E. Kranakis, D. Krizanc, and O. Morales-Ponce}

\institute{Loyola Marymount University, Los Angeles, USA \\ \email{jared.coleman@lmu.edu} \and
Carleton University, Ottawa, Canada \\ \email{kranakis@scs.carleton.ca} \and
Wesleyan University, Middletown, USA \\ \email{dkrizanc@wesleyan.edu} \and
California State University, Long Beach, USA \\ \email{Oscar.MoralesPonce@csulb.edu}}

\begin{document}
\maketitle

\begin{abstract}
We study online search on the real line by a unit-speed flying searcher for a target that starts at unknown distance $D \geq 1$ from the origin and moves directly away from the origin at constant speed $v < 1$. Unlike classical linear search, the searcher cannot observe the target while flying: detection is possible only at a landing, and only when the target lies between the origin and the landing point. We distinguish \emph{discovery}, where the process stops when the target is detected, from \emph{capture}, where the searcher must subsequently meet the target. When $v$ is known we give tight bounds for both objectives. When the initial distance is known but $v$ is not, a doubly-exponential landing strategy achieves competitive ratio $O((1-v)^{-23/4})$ for both objectives as $v \to 1^-$, and we conjecture that this exponent is optimal.
\keywords{Linear search \and escaping target \and competitive ratio \and online algorithms \and mobile agents}
\end{abstract}

\JC{This is a skeleton assembled from the two drafts in \texttt{references/}. Sections marked TODO still need to be written or verified; see \texttt{notes/model\_and\_results.md} for the current status of each claim.}

\FloatBarrier
\section{Introduction}\label{sec:intro}

Online search asks how efficiently a mobile agent can locate an object whose relevant parameters are not known in advance. The classical one-dimensional problem, often called linear search or the cow-path problem~\cite{beck1964linear,baezayates1993searching}, places a stationary target at an unknown point of the line and compares an online zig-zag trajectory with a clairvoyant searcher. The optimal deterministic competitive ratio is $9$, achieved by doubling the excursion length on alternating sides.

The present paper studies a mobile analogue with two additional constraints. First, the target is \emph{escaping}: it moves away from the origin at a fixed speed $v < 1$. Second, the searcher is a flying agent with \emph{landing-only} visibility. It receives no information while in flight. A target can be discovered only at a landing, and only if it lies on the segment joining the origin to the landing point. In this sense the searcher is ``forward blind'': overshooting the target in flight does not reveal it.

We distinguish two objectives. In the \emph{discovery} problem the cost is the time of the first landing that reveals the target. In the \emph{capture} problem the searcher must, after discovery, physically meet the target.

\subsection{Related work}\label{sec:related}

Demaine, Fekete, and Gal studied linear and star search with a fixed turn cost~\cite{demaine2006online}; Bose and De Carufel developed a general framework for searching on a line~\cite{bose2013revisiting,bose2017general}. Search for an oblivious target moving away from the origin was introduced in~\cite{coleman2022oblivious}, and the effect of knowing the target's speed or initial distance was studied in~\cite{coleman2024escaping,coleman2026knowledge}. The information regimes considered there motivate the four-way classification used here. Our model differs in a crucial sensing assumption: passing the target in flight does not reveal it. The landing-only restriction is also the one-dimensional analogue of the planar search-and-rescue model of~\cite{coleman2026searchrescueplane}, where an agent must land on the $x$-axis before it can search along it.

\JC{TODO: expand related work; cite the search-games books~\cite{alpern2003theory,gal1980} and the search-and-rescue line~\cite{coleman2023sar}.}

\subsection{Contributions and organization}

Section~\ref{sec:model} gives the model. Section~\ref{sec:static} treats the stationary target ($v = 0$). Section~\ref{sec:known-speed} analyzes geometric strategies when $v$ is known and $D$ is unknown, and Section~\ref{sec:lower} proves the matching lower bounds via Gal's theorem on sequence functionals. Section~\ref{sec:unknown-speed} treats the case where $D$ is known but $v$ is not. Section~\ref{sec:lower-conj} records a conjectured matching lower bound exponent for the latter case. Table~\ref{tab:summary} summarizes the results.

\begin{table}[t]
\centering
\caption{Summary of results. ``$D$ known'' means the magnitude $|d| = D$ is known but the sign of $d$ is not. Lower bounds marked $\ddagger$ are proved for alternating strategies (Section~\ref{sec:lower}); for arbitrary side patterns only the weaker bounds of Theorem~\ref{thm:lb-universal} are known.}
\label{tab:summary}
\begin{tabular}{llll}
\toprule
Objective & Information & Competitive ratio & Status \\
\midrule
Discovery & $v = 0$, $D$ unknown & $\frac{11 + 5\sqrt{5}}{2} \approx 11.09$ & tight$^\ddagger$ \\
Capture   & $v = 0$, $D$ unknown & $\frac{25}{2}$ & tight$^\ddagger$ \\
Discovery & $v$ known, $D$ unknown & Eq.~\eqref{eq:cr-disc-opt} & tight$^\ddagger$ \\
Capture   & $v$ known, $D$ unknown & Eq.~\eqref{eq:cr-cap-opt} & tight$^\ddagger$ \\
Discovery & $D$ known, $v$ unknown & $O\big((1-v)^{-23/4}\big)$ & upper bound \\
Capture   & $D$ known, $v$ unknown & $O\big((1-v)^{-23/4}\big)$ & upper bound \\
\bottomrule
\end{tabular}
\end{table}

\FloatBarrier
\section{Model and notation}\label{sec:model}

A searcher starts at the origin of $\mathbb{R}$ and moves with maximum speed $1$. A target starts at
\[
  d = \sigma D, \qquad \sigma \in \{-1, +1\}, \quad D \geq 1,
\]
and moves directly away from the origin at a constant speed $0 \leq v < 1$, so its position at time $t$ is $y(t) = \sigma(D + vt)$.

The searcher cannot detect the target while flying. At a landing at position $x$ at time $t$, it discovers the target if $x$ has the correct sign and $|x| \geq D + vt$; equivalently, the target lies between the origin and the landing point.

\paragraph{Sorties.} The landing-only restriction is meaningful only together with a restriction on how often the searcher may land. A searcher allowed to land at every point of its flight path detects the target at the moment it passes over it, and the problem collapses to classical linear search for an escaping target, whose optimal competitive ratio is $1 + 8(1+v)/(1-v)^2$ (and $9$ at $v = 0$)~\cite{alpern2003theory,coleman2022oblivious}. We therefore require that every flight start and end at the origin: the searcher's trajectory is a sequence of \emph{sorties}, each flying from the origin to a landing point and back. A landing that discovers the target ends the discovery problem; in the capture problem the searcher instead flies from the landing point directly to the target. One may think of a drone that must return to its base between flights.

\JC{This is the key modeling assumption, and neither draft states it. Both analyze sorties that return to the origin, so I believe it is what everyone has in mind, but we should agree on it explicitly. Without it the lower bounds below are false (classical search gives $9 < (11+5\sqrt5)/2$ at $v = 0$).}

A \emph{strategy} is a sequence of sortie lengths $a_i > 0$ and sides $\sigma_i \in \{-1, +1\}$, $i \geq 0$, landing at $x_i = \sigma_i a_i$. It is \emph{alternating} if $\sigma_i = (-1)^i$. All strategies analyzed in Sections~\ref{sec:static}--\ref{sec:unknown-speed} are alternating, and so is the class covered by the tight lower bounds of Section~\ref{sec:lower}; general side patterns are discussed in Section~\ref{sec:lower-general}. For an alternating strategy write
\[
  S_i := \sum_{j=0}^{i} a_j, \qquad S_{-1} := 0 .
\]
Because the searcher returns through the origin between consecutive landings, the $i$th landing occurs at time
\begin{equation}\label{eq:landing-time}
  t_i = 2 S_{i-1} + a_i = 2 S_i - a_i .
\end{equation}
The \emph{effective reach} of landing $i$ against speed $v$ is
\begin{equation}\label{eq:reach}
  B_i(v) := a_i - v t_i = (1-v) a_i - 2 v S_{i-1} ,
\end{equation}
and a target on the side of landing $i$ is discovered there exactly when $D \leq B_i(v)$. Equations~\eqref{eq:landing-time} and~\eqref{eq:reach} hold verbatim for non-alternating strategies, since landing times and reaches depend only on the sortie lengths.

Call sortie $i$ \emph{useful} if $B_i(v) \geq 1$ and $B_i(v) > B_j(v)$ for every earlier sortie $j$ on the same side. A sortie that is not useful discovers no target that was not already discovered, and deleting it shortens every later landing time, which by~\eqref{eq:reach} increases every later reach and decreases every discovery and capture time. Deleting a sortie from an alternating strategy, however, produces a non-alternating one. We say an alternating strategy is \emph{eventually monotone} if there is an index $i_0$ such that for all $i \geq i_0$, $B_i(v) \geq 1$ and $B_i(v) \geq B_j(v)$ for every $j < i$ with $j \equiv i \pmod 2$. Geometric strategies $a_i = c\, r^i$ with $r > (1+v)/(1-v)$ and $c \geq 1/(1-v)$ are eventually monotone with $i_0 = 0$.

If $\sigma$, $D$, and $v$ are all known, a unit-speed searcher meets the target at time
\begin{equation}\label{eq:opt}
  \opt(D, v) = \frac{D}{1 - v},
\end{equation}
and a clairvoyant searcher can land exactly at the meeting point, so we use~\eqref{eq:opt} as the benchmark for both discovery and capture. For known $v$ and unknown $D$ the competitive ratio of a strategy $\mathcal{A}$ is $\sup_{\sigma, D} T_{\mathcal{A}}(\sigma D, v) / \opt(D, v)$. When $D$ is known but $v$ is not, we use the speed-dependent ratio
\[
  \CR_{\mathcal{A}}(v) := \sup_{\sigma \in \{-1,+1\}} \frac{T_{\mathcal{A}}(\sigma D, v)}{D/(1-v)} ,
\]
and by scaling we may normalize $D = 1$.

Suppose the target is discovered at landing $i$. The searcher is then beyond the target and reverses direction, closing at relative speed $1 + v$, so the total capture time is
\begin{equation}\label{eq:capture-time}
  T_i^{\capt}(D, v) = t_i + \frac{a_i - D - v t_i}{1 + v} = \frac{2 S_i - D}{1 + v} .
\end{equation}

\FloatBarrier
\section{Stationary target}\label{sec:static}

\JC{TODO. Both drafts cover $v = 0$: geometric excursions $a_i = r^i$ give discovery ratio $r^2(r+1)/(r-1)$, minimized at $r = (1+\sqrt5)/2$ with value $(11+5\sqrt5)/2$, and capture ratio $2r^3/(r-1) - 1$, minimized at $r = 3/2$ with value $25/2$. The lower bounds are the $v = 0$ case of Theorems~\ref{thm:lb-disc} and~\ref{thm:lb-cap}; the monotone-sequence technique of EK's notes (\texttt{references/dis-cap-line.pdf}, Sections~2.2.2 and 2.2.3) does not close because the discovery and capture functionals involve three consecutive sorties, which is exactly what Gal's theorem handles. Decide whether this section should exist at all or whether $v = 0$ is just a corollary.}

\FloatBarrier
\section{Known speed, unknown distance}\label{sec:known-speed}

Assume $v$ is known, whereas $D \geq 1$ and $\sigma$ are unknown. Consider geometric excursions $a_i = r^i$ with $r > 1$, so $S_i = (r^{i+1} - 1)/(r - 1)$. The effective reach~\eqref{eq:reach} is eventually positive only if
\begin{equation}\label{eq:feasible}
  r > \frac{1 + v}{1 - v},
\end{equation}
since $B_i(v)/a_i \to ((1-v) r - (1+v))/(r-1)$.

\subsection{Discovery}

A target on the side of landing $i$ that is first discoverable there was missed at landing $i-2$, so the worst case is $D \downarrow B_{i-2}(v)$ with discovery cost $t_i$. Asymptotically in $i$,
\begin{equation}\label{eq:cr-disc-r}
  \CR_{\disc}(r, v) = (1 - v) \frac{r^2 (r + 1)}{(1 - v) r - (1 + v)} .
\end{equation}

\begin{theorem}[Known-speed discovery]\label{thm:disc-known}
For landing-only discovery with unknown initial distance and known speed $v$, the optimal geometric expansion is
\begin{equation}\label{eq:r-disc-opt}
  r^{*}_{\disc}(v) = \frac{1 + 2v + \sqrt{5 + 4v}}{2 (1 - v)},
\end{equation}
and the resulting competitive ratio is
\begin{equation}\label{eq:cr-disc-opt}
  \CR^{*}_{\disc}(v) = \frac{\big(\sqrt{5 + 4v} + 3\big)\big(2v + \sqrt{5 + 4v} + 1\big)^2}{4 (1 - v)^2 \big(\sqrt{5 + 4v} - 1\big)} .
\end{equation}
At $v = 0$ this is $(11 + 5\sqrt5)/2$, and $\CR^{*}_{\disc}(v) \sim 27/(1-v)^2$ as $v \to 1^-$.
\end{theorem}
\begin{proof}
Minimize $f(r) = r^2(r+1)/((1-v) r - (1+v))$ over the feasible range~\eqref{eq:feasible}. Then $f'(r) = 0$ if and only if $(1-v) r^2 - (1 + 2v) r - (1+v) = 0$, whose feasible root is~\eqref{eq:r-disc-opt}; substituting into~\eqref{eq:cr-disc-r} gives~\eqref{eq:cr-disc-opt}. These identities are checked symbolically in \texttt{code/verify\_known\_speed.wls}.
\end{proof}

The matching lower bound is Theorem~\ref{thm:lb-disc} in Section~\ref{sec:lower}.

\subsection{Capture}

With the same worst case $D \downarrow B_{i-2}(v)$, equation~\eqref{eq:capture-time} gives
\begin{equation}\label{eq:cr-cap-r}
  \CR_{\capt}(r, v) = \frac{1 - v}{1 + v} \left( \frac{2 r^3}{(1 - v) r - (1 + v)} - 1 \right) .
\end{equation}

\begin{theorem}[Known-speed capture]\label{thm:cap-known}
The optimal geometric expansion for capture is
\begin{equation}\label{eq:r-cap-opt}
  r^{*}_{\capt}(v) = \frac{3 (1 + v)}{2 (1 - v)},
\end{equation}
and the resulting competitive ratio is
\begin{equation}\label{eq:cr-cap-opt}
  \CR^{*}_{\capt}(v) = \frac{27 (1 + v)}{2 (1 - v)^2} - \frac{1 - v}{1 + v} = \frac{(2v + 1)(v + 5)^2}{2 (1 - v)^2 (1 + v)} .
\end{equation}
At $v = 0$ this is $25/2$, and $\CR^{*}_{\capt}(v) \sim 27/(1-v)^2$ as $v \to 1^-$.
\end{theorem}
\begin{proof}
It suffices to minimize $K(r) = r^3 / ((1-v) r - (1+v))$. The derivative vanishes when $3 r^2 ((1-v) r - (1+v)) = (1-v) r^3$, i.e.\ $2(1-v) r = 3(1+v)$, giving~\eqref{eq:r-cap-opt} and $K^{*} = 27(1+v)^2 / (4(1-v)^3)$. Substituting into~\eqref{eq:cr-cap-r} yields~\eqref{eq:cr-cap-opt}. See \texttt{code/verify\_known\_speed.wls}.
\end{proof}

The matching lower bound is Theorem~\ref{thm:lb-cap}.

\EK{The formulas for $\CR_{\disc}(v)$ and $\CR_{\capt}(v)$ in \texttt{references/dis-cap-line.pdf} (Table~1 and Sections~2.3.2--2.3.3) differ from~\eqref{eq:cr-disc-opt} and~\eqref{eq:cr-cap-opt} for $v > 0$ although they agree at $v = 0$. The simulator \texttt{code/simulate.py} agrees with the formulas above; see \texttt{notes/model\_and\_results.md}.}

\FloatBarrier
\section{Lower bounds for known speed}\label{sec:lower}

Throughout this section $v \in [0, 1)$ is fixed and $q := (1+v)/(1-v)$. The lower bounds follow the standard route for linear search: an adversary places the target just beyond the reach of a landing, the resulting ratios form a sequence of functionals of the strategy, and a theorem of Gal shows that the supremum of such functionals is minimized by a geometric sequence. The one twist is that the reaches $B_i(v)$ are not positive linear forms in the sortie lengths, so we parametrize strategies by their reaches instead.

\subsection{Reach parametrization}

\begin{lemma}[Reach parametrization]\label{lem:reach}
Let $(a_i)_{i \geq 0}$ be an alternating strategy and $z_i := B_i(v)$ its reaches. Then
\begin{equation}\label{eq:S-from-z}
  S_i = q\, S_{i-1} + \frac{z_i}{1-v} = \frac{1}{1-v} \sum_{j=0}^{i} q^{\,i-j} z_j, \qquad
  a_i = \frac{z_i + 2v S_{i-1}}{1-v}, \qquad
  t_i = S_i + S_{i-1} .
\end{equation}
Conversely, every sequence of positive reals $(z_i)_{i \geq 0}$ is the reach sequence of exactly one alternating strategy, and that strategy has positive sortie lengths.
\end{lemma}
\begin{proof}
By~\eqref{eq:reach}, $z_i = (1-v) a_i - 2v S_{i-1} = (1-v) S_i - (1+v) S_{i-1}$, which rearranges to the recurrence in~\eqref{eq:S-from-z}; unrolling it with $S_{-1} = 0$ gives the sum. The formulas for $a_i = S_i - S_{i-1}$ and $t_i = 2S_{i-1} + a_i$ follow. For the converse, define $S_i$ by the recurrence and $a_i := S_i - S_{i-1} = \frac{2v}{1-v} S_{i-1} + \frac{z_i}{1-v} > 0$. The identities are checked symbolically in \texttt{code/verify\_lower\_bound.wls}.
\end{proof}

For a positive sequence $Z = (z_0, z_1, \ldots)$ we write $S_i(Z)$ for the quantity defined by~\eqref{eq:S-from-z}; it is a linear form in $z_0, \ldots, z_i$ with positive coefficients that depend only on the differences $i - j$.

\subsection{Gal's theorem}

For a sequence $X = (x_0, x_1, \ldots)$ let $X^{+1} := (x_1, x_2, \ldots)$ and let $G_\alpha := (1, \alpha, \alpha^2, \ldots)$.

\begin{theorem}[Gal~\cite{gal1980}; Schuierer~\cite{schuierer2001lower}]\label{thm:gal}
Let $X = (x_0, x_1, \ldots)$ be a sequence of positive numbers, $r \geq 0$ an integer, and $\alpha := \limsup_{n \to \infty} x_n^{1/n} \in \mathbb{R} \cup \{+\infty\}$. Let $F_i$, $i \geq 0$, be functionals such that
\begin{enumerate}
  \item $F_i(X)$ depends only on $x_0, \ldots, x_{i+r}$;
  \item $F_i$ is continuous on $\{x_k > 0 : 0 \leq k \leq i + r\}$;
  \item $F_i(\lambda X) = F_i(X)$ for all $\lambda > 0$;
  \item $F_i(X + Y) \leq \max\{F_i(X), F_i(Y)\}$;
  \item $F_{i+1}(X) \geq F_i(X^{+1})$.
\end{enumerate}
Then $\sup_{i \geq 0} F_i(X) \geq \sup_{i \geq 0} F_i(G_\alpha)$, where for $\alpha = +\infty$ the right-hand side is read as $\lim_{\alpha \to \infty} \sup_i F_i(G_\alpha)$.
\end{theorem}

\JC{Statement transcribed from Theorem~3 of~\cite{angelopoulos2019bestoftwo} (there condition~5 is printed as $F_{i+1}(X) \geq F_i(X^{k+1})$, surely a typo for $X^{+1}$) and from Schuierer's definition $X^{+k} = (x_k, x_{k+1}, \ldots)$. Please check against~\cite{schuierer2001lower} and Alpern--Gal~\cite{alpern2003theory} before submission, in particular the convention for $\alpha = \infty$.}

The functionals we use are all of the form $F_i(Z) = c_1 \frac{L_i(Z)}{z_i} - c_0$ with $c_1 > 0$, $c_0 \geq 0$, and $L_i$ a positive linear form in $z_0, \ldots, z_{i+2}$ whose coefficient of $z_j$ depends only on $i - j$. Such functionals satisfy the hypotheses of Theorem~\ref{thm:gal} with $r = 2$: conditions~1--3 are immediate; condition~4 is the mediant inequality $\frac{p_1 + p_2}{q_1 + q_2} \leq \max\{\frac{p_1}{q_1}, \frac{p_2}{q_2}\}$ for positive $p_k, q_k$, which is unaffected by the affine map $u \mapsto c_1 u - c_0$; and condition~5 holds because $F_{i+1}(Z) - F_i(Z^{+1})$ equals $c_1$ times the coefficient of $z_0$ in $L_{i+1}(Z)$, times $z_0 / z_{i+1} \geq 0$.

For the geometric sequence $G_\alpha$ with $\alpha \neq q$,~\eqref{eq:S-from-z} gives
\begin{equation}\label{eq:S-geo}
  S_i(G_\alpha) = \frac{\alpha^{i+1} - q^{i+1}}{(1-v)(\alpha - q)}, \qquad\text{and}\qquad S_i(G_q) = \frac{(i+1) q^i}{1-v}.
\end{equation}

\subsection{Discovery}

\begin{theorem}[Known-speed discovery, lower bound]\label{thm:lb-disc}
Every eventually monotone alternating strategy has discovery competitive ratio at least $\CR^{*}_{\disc}(v)$, given by~\eqref{eq:cr-disc-opt}.
\end{theorem}
\begin{proof}
Let $Z = (z_i)$ be the reach sequence of the strategy and $i_0$ the index from the definition of eventual monotonicity. Fix $i \geq i_0$ and $\varepsilon > 0$, and place the target on side $(-1)^i$ at distance $D = z_i + \varepsilon \geq 1$. No landing $j \leq i$ on that side discovers it, because $B_j(v) \leq z_i < D$. The next landing on that side is landing $i+2$, so the discovery time is at least $t_{i+2} = S_{i+2} + S_{i+1}$ and
\[
  \CR_{\disc} \;\geq\; (1-v)\,\frac{S_{i+2} + S_{i+1}}{z_i + \varepsilon} \quad\text{for all } \varepsilon > 0,
  \qquad\text{hence}\qquad
  \CR_{\disc} \;\geq\; F_i(Z) := (1-v)\,\frac{S_{i+2}(Z) + S_{i+1}(Z)}{z_i} .
\]
By Lemma~\ref{lem:reach}, $F_i$ has the form discussed after Theorem~\ref{thm:gal}, so Theorem~\ref{thm:gal} applies. Let $Z' := (z_{i_0}, z_{i_0+1}, \ldots)$, a sequence of positive numbers. Iterating condition~5, $F_i(Z') \leq F_{i + i_0}(Z) \leq \CR_{\disc}$ for every $i \geq 0$, so with $\alpha := \limsup_n (z'_n)^{1/n} \in [1, +\infty]$,
\[
  \CR_{\disc} \;\geq\; \sup_{i} F_i(Z') \;\geq\; \sup_i F_i(G_\alpha) .
\]
It remains to evaluate the right-hand side using~\eqref{eq:S-geo}. If $\alpha > q$ then
\[
  F_i(G_\alpha) \;\xrightarrow[i \to \infty]{}\; (1-v)\,\frac{\alpha^2 (\alpha + 1)}{(1-v)\alpha - (1+v)} = \CR_{\disc}(\alpha, v),
\]
the geometric-strategy ratio~\eqref{eq:cr-disc-r}, and $\sup_i F_i(G_\alpha)$ is at least this limit. If $1 \leq \alpha \leq q$ then $F_i(G_\alpha) \to \infty$, since the $q^{i}$ terms dominate (for $\alpha = q$, $F_i(G_q)$ grows linearly in $i$), and if $\alpha = \infty$ the bound is trivial. In every case $\CR_{\disc} \geq \inf_{\alpha > q} \CR_{\disc}(\alpha, v) = \CR^{*}_{\disc}(v)$, where the last equality is the minimization carried out in the proof of Theorem~\ref{thm:disc-known}: $\CR_{\disc}(\cdot, v)$ tends to $+\infty$ at both ends of $(q, \infty)$ and has the unique critical point $r^{*}_{\disc}(v)$ there. The limits and the minimization are verified in \texttt{code/verify\_lower\_bound.wls}.
\end{proof}

\subsection{Capture}

\begin{theorem}[Known-speed capture, lower bound]\label{thm:lb-cap}
Every eventually monotone alternating strategy has capture competitive ratio at least $\CR^{*}_{\capt}(v)$, given by~\eqref{eq:cr-cap-opt}.
\end{theorem}
\begin{proof}
With the same adversary as in Theorem~\ref{thm:lb-disc}, the target $D = z_i + \varepsilon$ on side $(-1)^i$ is discovered at some landing $j \geq i+2$, and by~\eqref{eq:capture-time} the capture time is $(2S_j - D)/(1+v) \geq (2S_{i+2} - D)/(1+v)$. Dividing by $D/(1-v)$ and letting $\varepsilon \to 0$,
\[
  \CR_{\capt} \;\geq\; F'_i(Z) := \frac{1-v}{1+v} \left( \frac{2 S_{i+2}(Z)}{z_i} - 1 \right) ,
\]
again a functional of the admissible form. Theorem~\ref{thm:gal} applied to $Z' = (z_{i_0}, z_{i_0+1}, \ldots)$ gives $\CR_{\capt} \geq \sup_i F'_i(G_\alpha)$, and by~\eqref{eq:S-geo}, for $\alpha > q$,
\[
  F'_i(G_\alpha) \;\xrightarrow[i \to \infty]{}\; \frac{1-v}{1+v}\left( \frac{2\alpha^3}{(1-v)\alpha - (1+v)} - 1 \right) = \CR_{\capt}(\alpha, v),
\]
while $F'_i(G_\alpha) \to \infty$ for $\alpha \leq q$. Hence $\CR_{\capt} \geq \inf_{\alpha > q} \CR_{\capt}(\alpha, v) = \CR^{*}_{\capt}(v)$ by the minimization in the proof of Theorem~\ref{thm:cap-known}.
\end{proof}

\begin{remark}[Exact tightness]\label{rem:tight}
The upper bounds of Theorems~\ref{thm:disc-known} and~\ref{thm:cap-known} were stated asymptotically in $i$. They hold exactly, as suprema over all $D \geq 1$, for $a_i = c\, r^{*} {}^i$ with $c = 1/(1-v)$: for a target just beyond $B_{i-2}(v)$ the ratio is
\[
  (1-v)\,\frac{c\,(r^{i+1} + r^{i}) - 2}{c\, r^{i-2}\big((1-v) r - (1+v)\big) + 2v} \;<\; \CR_{\disc}(r, v)
\]
for every finite $i \geq 2$, and the targets $D \in [1, B_0(v)]$ and $D \in [1, B_1(v)]$ found at the first two landings have ratios at most $(1-v) t_1 = 2 + r^{*}$, which is below $\CR^{*}$. The capture case is identical. The script \texttt{code/explore\_strategies.py} confirms numerically that the full supremum equals the asymptotic value. \JC{Write this out properly or drop the remark.}
\end{remark}

\subsection{General side patterns}\label{sec:lower-general}

Theorems~\ref{thm:lb-disc} and~\ref{thm:lb-cap} cover alternating strategies. Shortening every useless sortie to length zero shows that an arbitrary strategy is equivalent to an alternating strategy in which some sorties have length zero and hence negative reach; the functionals $F_i$ and $F'_i$ are undefined at such indices, and the natural repair (dividing by the largest reach so far on the same side instead of $z_i$) destroys condition~4 of Theorem~\ref{thm:gal}, because a maximum of linear forms is not superadditive. Whether non-alternating side patterns can help is therefore open. Two pieces of evidence suggest they cannot.

First, numerical optimization over periodic side patterns such as $(+,+,-)$, $(+,+,-,-)$ and $(+,+,-,+,-,-)$, with free sortie lengths within the period and a free period ratio (\texttt{code/explore\_strategies.py}), always returns the alternating optimum, with the surplus sorties shrinking to zero, for both objectives and for $v \in \{0, 0.5\}$.

Second, the adversary argument does give a weaker bound for every strategy. For an arbitrary strategy whose sorties are all useful, the target just beyond $z_i$ is discovered no earlier than the \emph{next} landing $i+1$, and Lemma~\ref{lem:reach} holds for any side pattern, so $\CR_{\disc} \geq (1-v)(S_{i+1} + S_i)/z_i$ and $\CR_{\capt} \geq \frac{1-v}{1+v}(2S_{i+1}/z_i - 1)$ for all $i$. Theorem~\ref{thm:gal} and~\eqref{eq:S-geo} then yield the following.

\begin{theorem}[Universal lower bounds]\label{thm:lb-universal}
Every strategy, with any side pattern, satisfies
\[
  \CR_{\disc} \;\geq\; \frac{3 + v + 2\sqrt{2(1+v)}}{1-v}
  \qquad\text{and}\qquad
  \CR_{\capt} \;\geq\; \frac{8}{1-v} - \frac{1-v}{1+v} .
\]
At $v = 0$ these are $3 + 2\sqrt2 \approx 5.83$ and $7$.
\end{theorem}
\begin{proof}
For $\alpha > q$ the two functionals on $G_\alpha$ tend to $(1-v)\alpha(\alpha+1)/((1-v)\alpha - (1+v))$ and $\frac{1-v}{1+v}\big(2\alpha^2/((1-v)\alpha - (1+v)) - 1\big)$, minimized over $\alpha > q$ at $\alpha = (1 + v + \sqrt{2(1+v)})/(1-v)$ and $\alpha = 2q$ respectively; for $\alpha \leq q$ they diverge. See \texttt{code/verify\_lower\_bound.wls}.
\end{proof}

\begin{conjecture}\label{conj:alternating}
Theorems~\ref{thm:lb-disc} and~\ref{thm:lb-cap} hold for every strategy, with any side pattern.
\end{conjecture}

\FloatBarrier
\section{Known distance, unknown speed}\label{sec:unknown-speed}

Assume $D = |d|$ is known (normalize $D = 1$) but $\sigma$ and $v$ are unknown. No uniform constant ratio is possible as $v \to 1^-$, so we bound the ratio as a function of the actual speed.

\subsection{Doubly-exponential strategy and induced speed thresholds}

Fix $A > 1$ and $b > 1$ and use alternating excursions
\begin{equation}\label{eq:doubly-exp}
  a_i = A^{b^i} .
\end{equation}
The maximal speed handled by round $i$ is defined by the critical equality $1 + v_i t_i = a_i$, so by~\eqref{eq:landing-time}
\begin{equation}\label{eq:threshold}
  v_i = \frac{a_i - 1}{a_i + 2 S_{i-1}}, \qquad \delta_i := 1 - v_i = \frac{1 + 2 S_{i-1}}{a_i + 2 S_{i-1}} .
\end{equation}
Since the sequence is doubly exponential, $S_i = a_i(1 + o(1))$ and $a_{i-1}/a_i = A^{-(b-1) b^{i-1}}$, hence
\begin{equation}\label{eq:delta-asymp}
  \delta_i \sim 2 \frac{a_{i-1}}{a_i} = 2 A^{-(b-1) b^{i-1}} .
\end{equation}
For a fixed actual speed $v$ let $i(v) := \max\{ i : v_i < v \}$. Round $i(v) + 1$ is speed-sufficient; since the sign is unknown, the adversary can force one more round, so discovery happens by round $i(v) + 2$.

\begin{theorem}[Unknown-speed discovery]\label{thm:disc-unknown}
For the strategy~\eqref{eq:doubly-exp},
\[
  \CR_{\disc}(v) = O\big( (1 - v)^{-\alpha(b)} \big), \qquad \alpha(b) := \frac{b^3}{b - 1} - 1,
\]
as $v \to 1^-$. The exponent is minimized uniquely at $b = 3/2$, giving $\CR_{\disc}(v) = O\big((1-v)^{-23/4}\big)$.
\end{theorem}
\begin{proof}
\JC{TODO: port the proof from the draft (Theorem~5.1 there). Sketch: with $i = i(v)$, $\delta_{i+1} \leq 1 - v < \delta_i$; the normalized cost $\delta\, t_{i+2}$ is increasing in $\delta$, so the worst case on this interval is $\delta_i a_{i+2} \asymp A^{[b^3 - (b-1)] b^{i-1}}$ while $\delta_i^{-1} \asymp A^{(b-1) b^{i-1}}$. Then $\alpha'(b) = b^2 (2b - 3)/(b-1)^2$.}
\end{proof}

\begin{theorem}[Unknown-speed capture]\label{thm:cap-unknown}
For the strategy~\eqref{eq:doubly-exp}, $\CR_{\capt}(v) = O\big((1 - v)^{-\alpha(b)}\big)$ with the same $\alpha(b)$, hence $O\big((1-v)^{-23/4}\big)$ at $b = 3/2$.
\end{theorem}
\begin{proof}
\JC{TODO. By~\eqref{eq:capture-time} with $D = 1$, $\CR_{\capt}(v) \leq \frac{1-v}{1+v}(2 S_{i(v)+2} - 1) \leq (1 + o(1)) (1-v) a_{i(v)+2}$, and the exponent calculation is identical to Theorem~\ref{thm:disc-unknown}.}
\end{proof}

\FloatBarrier
\section{A conjectured matching lower exponent}\label{sec:lower-conj}

Normalize $D = 1$ and consider an arbitrary alternating landing strategy with $\delta_i = (1 + 2 S_{i-1})/(S_i + S_{i-1})$ as in~\eqref{eq:threshold}. When $S_i \gg S_{i-1}$, $\delta_i \asymp S_{i-1}/S_i$. Taking a speed just above $v_i$ and the adverse sign, both discovery and capture cost roughly $\delta_i S_{i+2}$ in normalized terms. With local growth exponents $p_i := \log S_i / \log S_{i-1}$, this suggests the local lower-bound exponent $p_i p_{i+1} p_{i+2}/(p_i - 1) - 1$, which for constant $p$ is $p^3/(p-1) - 1$, minimized at $p = 3/2$ with value $23/4$.

\begin{conjecture}[Universal unknown-speed lower exponent]\label{conj:lower}
For every deterministic landing-only strategy $\mathcal{A}$ there is a sequence $v_k \to 1^-$ such that, for discovery and likewise for capture,
\[
  \CR_{\mathcal{A}}(v_k) = \Omega\big( (1 - v_k)^{-23/4 + o(1)} \big) .
\]
\end{conjecture}

\FloatBarrier
\section{Conclusion and open problems}\label{sec:conclusion}

\JC{TODO. Open directions from the drafts: (1) prove or refute Conjecture~\ref{conj:lower}; (2) extend the known-speed lower bounds from alternating strategies to arbitrary side patterns (Conjecture~\ref{conj:alternating}); (3) nonstationary doubly-exponential schedules (varying $b$) to improve lower-order factors; (4) randomized strategies; (5) positive landing times or turn costs, connecting to~\cite{demaine2006online}.}

\bibliographystyle{splncs04}
\bibliography{refs}

\end{document}
